\documentclass[../../main.tex]{subfiles}
\graphicspath{{../../image/}{../../figure/}}
\begin{document}

\section{特征线的传播与 $N$ 波的发展}\label{section:ch11-特征线传播与N波发展}

真正非线性的另一影响解大时间行为的特征是特征线的传播。为看清这一机制的效应, 我们研究初值具紧支集时解的渐近行为。由\cref{theorem:ch11-11-5-2}我们已知, 振幅以 $O(t^{-1/2})$ 衰减。这里更细致的考察将揭示, 渐近地, 解达到 $N$ 波剖面, 即两侧由激波夹持的中心稀疏波, 激波振幅以 $O(t^{-1/2})$ 衰减。

\begin{theorem}\label{theorem:ch11-11-6-1}
设 $u$ 是初值为 $u_0$ 的可容许解, $u_0(x)=0$ 对 $|x|>\ell$。考虑从 $(-\ell,0)$ 出发的极小前向特征线 $\chi_-(\cdot)$ 和从 $(\ell,0)$ 出发的极大前向特征线 $\chi_+(\cdot)$。则
\begin{align}
u(x\pm,t)=0,\quad \text{对 }t>0\text{ 且 }x<\chi_-(t)\text{ 或 }x>\chi_+(t).\label{align:ch11-11-6-1}
\end{align}
当 $t\to\infty$ 时,
\begin{align}
f'(u(x\pm,t))=\frac{x}{t}+O\left(\frac1t\right),\quad \chi_-(t)<x<\chi_+(t),\label{align:ch11-11-6-2}\\
u(x\pm,t)=\frac{1}{f''(0)}\left[\frac{x}{t}-f'(0)\right]+O\left(\frac1t\right),\quad \chi_-(t)<x<\chi_+(t),\label{align:ch11-11-6-3}\\
\begin{cases}
\chi_-(t)=tf'(0)-[2q_-tf''(0)]^{1/2}+O(1),\\
\chi_+(t)=tf'(0)+[2q_+tf''(0)]^{1/2}+O(1),
\end{cases}\label{align:ch11-11-6-4}
\end{align}
其中 $q_-$ 与 $q_+$ 由\eqref{align:ch11-11-4-4}给出。此外, $u(\cdot,t)$ 在区间 $[\chi_-(t),\chi_+(t)]$ 上的递减变差为 $O(t^{-1})$。
\end{theorem}

\begin{proof}
由于 $\chi_-(\cdot)$ 极小而 $\chi_+(\cdot)$ 极大, 任意点 $(x,t)$, $t>0$, $x<\chi_-(t)$ 或 $x>\chi_+(t)$, 出发的极值后向特征线将在 $u_0$ 的支集之外被 $x$ 轴截获。这证明了\eqref{align:ch11-11-6-1}。

另一方面, 从满足 $t>0$ 且 $\chi_-(t)<x<\chi_+(t)$ 的点 $(x,t)$ 出发的极小或极大后向特征线 $\xi(\cdot)$, 将在区间 $[-\ell,\ell]$ 内被 $x$ 轴截获, 即 $\xi(0)\in[-\ell,\ell]$。因此当 $t\to\infty$ 时, $x-tf'(u(x\pm,t))=\xi(0)=O(1)$, 这给出\eqref{align:ch11-11-6-2}。\footnote{当 $t\to\infty$ 时, $\xi(0)$ 聚集于分界线发源点集。在\eqref{align:ch11-11-4-2}以 $\bar u=0$ 仅在单点 $\bar x$ 处成立的一般情形, 将该点归一化为 $\bar x=0$, 则 $\xi(0)$ 聚集于原点, 因而\eqref{align:ch11-11-6-2}中的 $O(t^{-1})$ 升级为 $o(t^{-1})$。此外, 若 $u_0$ 是 $C^1$ 且 $u_0'(0)>0$, 则\eqref{align:ch11-11-6-2}中的 $O(t^{-1})$ 改进为 $O(t^{-2})$, 且 $t$ 充分大时, $u(\cdot,t)$ 在区间 $(\chi_-(t),\chi_+(t))$ 上是 $C^1$ 的。}

由\cref{theorem:ch11-11-5-2}, 当 $t\to\infty$ 时 $u=O(t^{-1/2})$, 因而假设 $f$ 是 $C^3$ 的, $f'(u)=f'(0)+f''(0)u+O(t^{-1})$。于是\eqref{align:ch11-11-6-3}由\eqref{align:ch11-11-6-2}推出。

为推导 $\chi_\pm(t)$ 当 $t\to\infty$ 的渐近式, 首先注意由 $0\geqslant\dot\chi_-(t)-f'(0)\geqslant O(t^{-1/2})$, $0\leqslant\dot\chi_+(t)-f'(0)\leqslant O(t^{-1/2})$, 这进而给出 $0\geqslant\chi_-(t)-tf'(0)\geqslant O(t^{1/2})$, $0\leqslant\chi_+(t)-tf'(0)\leqslant O(t^{1/2})$。接下来援引\cref{theorem:ch11-11-4-2}: 分界线 $x=\bar x+tf'(0)$ 发自点 $(\bar x,0)$, $\bar x\in[-\ell,\ell]$, 沿其上 $u$ 为零, 且对任意 $t>0$,
\begin{align}
\int_{\bar x+tf'(0)}^{\chi_-(t)}u(y,t)\,\mathrm{d}y=q_-,\qquad
\int_{\bar x+tf'(0)}^{\chi_+(t)}u(y,t)\,\mathrm{d}y=q_+.\label{align:ch11-11-6-5}
\end{align}
在\eqref{align:ch11-11-6-5}中代入 $u$ 的渐近式\eqref{align:ch11-11-6-3}, 经简单积分推出
\begin{align}
\frac{1}{2q_\pm tf''(0)}[\chi_\pm(t)-tf'(0)]^2=1+O(t^{-1/2}),\label{align:ch11-11-6-6}
\end{align}
由此得\eqref{align:ch11-11-6-4}。证毕。
\end{proof}

\end{document}
